\section{Deployment des Frameworks}
Das Framework wird auf einem zentralen Server (siehe Tabelle \ref{tab:jupyterhub_vm}) des Instituts für die Digitalisierung von Arbeits- und Lebenswelten (IDiAL) der Fachhochschule Dortmund gehostet.
Ziel ist es, einen benutzerdefinierten Python-Kernel bereitzustellen,
der allen Benutzern zentral zur Verfügung steht und auf das Framework samt Abhängigkeiten zugreifen kann.
Die Paketverwaltung erfolgt mithilfe von Poetry Version 2.1.2, um eine reproduzierbare und isolierte Entwicklungsumgebung sicherzustellen.
Das Zielsystem stellte standardmäßig Python in Version 3.8.10 bereit,
welche jedoch nicht kompatibel mit allen Abhängigkeiten des Projekts war.
Daher wurde Python 3.10 von python.org heruntergeladenen und manuell im Home-Verzeichnisses des Benutzers philipp kompiliert und installiert.
Dabei wurde die Systeminstallation Python 3.8.10 nicht verändern oder gelöscht:

\begin{center}
    \begin{verbatim}
        ./configure --prefix=$HOME/python-3.10
        make
        make install
        export PATH="$HOME/python-3.10/bin:$PATH"
        source ~/.bashrc
    \end{verbatim}
\end{center}

Anschließend wurde das Repositorie von GitHub geklont.
Da es sich um ein privates Repositorie handelt, wurde ein SSH-Schlüssel generiert und dieser bei GitHub eingerichtet.
Danach wurde das Projekt eingerichtet und anschließénd in das Verzeichnis /opt verschoben um eine zentrale Nutzung zu ermöglichen:

\begin{center}
    \begin{verbatim}
        git clone git@github.com:PSoldOut/visualkinematics_v2.git
        cd visualkinematics_v2
        poetry install
        sudo mv ../visualkinematics_v2 /opt
    \end{verbatim}
\end{center}

Die dadurch erstellte virtuelle Umgebung wurde als globaler Kernel für alle JupyterHub-Nutzer registriert. Dazu wurde
folgender Befehl genutzt:



 \begin{center}
    \begin{verbatim}
        'sudo $(poetry run which python) -m ipykernel install
        --name=visualkinematics_v2
        --display-name "VisualKinematics v2"
        --prefix=/usr/local'
    \end{verbatim}
\end{center}

 Dies installiert den Kernel systemweit in /usr/local/share/jupyter/kernels/.
 Danach ist der Kernel im JupyterHub sichtbar und kann von allen Nutzern verwendet werden. Der Zugriff auf
 das Poetry-Projekt und dessen Bibliotheken ist gegeben. Nutzer können nun direkt mit dem
 Framework in Jupyter arbeiten.
 Damit die Schieberegler und Knöpfe des Frameworks ordnungsgemäß Funktionieren,
 musste ipyevents und die zugehörige Labextension auf dem Host installiert werden:

 \begin{center}
    \begin{verbatim}
        sudo /opt/tljh/user/bin/python3 -m pip install ipyevents
        sudo /opt/tljh/user/bin/jupyter labextension install ipyevents
    \end{verbatim}
\end{center}